\section{AI Transparency Log}
\label{app:ai_log}

\subsection*{Tools Used}
\begin{itemize}[leftmargin=1.5em, itemsep=2pt]
    \item \textbf{ChatGPT (OpenAI) and Gemini (Google):} Used as interactive
    thought partners for document structuring, LaTeX troubleshooting, phrasing
    refinement, and initial consistency audits across sections.
\end{itemize}

\subsection*{AI-Assisted Work and Deep Claim Audits}
Generative AI suggestions were treated as preliminary drafts and independently
audited. Four representative audits that materially affected mathematical or
numerical claims are recorded below.

\paragraph{Audit 1: Step-doubling controller resizing exponent
(Section~2.2.4 and Meichen Ding ICS).}
\begin{itemize}[leftmargin=1.5em, itemsep=2pt]
    \item \textbf{Task/query:} Formulate the step-size update based on the
    step-doubling Richardson error estimate $E$.
    \item \textbf{Raw AI output:} $h_{\mathrm{new}}=hE^{-1/p}$, where $p$ is
    the global order of the base integrator.
    \item \textbf{Risk audited:} This conflates global error
    $\mathcal{O}(h^p)$ with a local defect of order $\mathcal{O}(h^{p+1})$.
    \item \textbf{Verification and action:} The controller was re-derived with
    exponent $-1/(p+1)$, safety factor $0.9$, and limiter $[0.2,2.0]$.
\end{itemize}

\paragraph{Audit 2: Frozen-Jacobian stability interpretation
(Section~2.3.4 and Jingyan Wei ICS).}
\begin{itemize}[leftmargin=1.5em, itemsep=2pt]
    \item \textbf{Task/query:} Assess explicit stability constraints for
    $\dot X=X(I_n-X^TX)$.
    \item \textbf{Raw AI output:} The condition
    $|1+h\lambda_j|\leq1$, giving $h<2/26\approx0.0769$, was presented as a
    sufficient certificate of stability over $[0,2]$.
    \item \textbf{Risk audited:} $J_f(X(t))$ changes with the nonlinear
    trajectory, so the initial frozen spectrum is only a local diagnostic.
    \item \textbf{Verification and action:} Section~2.3.4 was revised to state
    this limitation, and Section~4 uses empirical step-size sweeps to test the
    nonlinear behaviour.
\end{itemize}

\paragraph{Audit 3: Newton tolerance versus time-discretisation error
(Section~2.4.2 and Tianrui Ji ICS).}
\begin{itemize}[leftmargin=1.5em, itemsep=2pt]
    \item \textbf{Task/query:} Determine whether a fixed Newton residual
    tolerance $10^{-12}$ guarantees the theoretical order of Implicit Euler.
    \item \textbf{Raw AI output:} A sufficiently small fixed tolerance was
    claimed to guarantee convergence without affecting global order.
    \item \textbf{Risk audited:} Algebraic and time-discretisation errors are
    different; asymptotically, a fixed residual can eventually dominate.
    \item \textbf{Verification and action:} The report restricts the claim to
    the tested finite grids. Tightening the tolerance to $10^{-14}$ changes the
    terminal state by about $1.14\times10^{-8}$, below the tested Implicit
    Euler discretisation error.
\end{itemize}

\paragraph{Audit 4: Stiff damping of Gauss IRK4
(Section~5.4 and the Jingyan Wei and Meichen Ding ICSs).}
\begin{itemize}[leftmargin=1.5em, itemsep=2pt]
    \item \textbf{Task/query:} Analyse the two-stage fourth-order
    Gauss--Legendre method on stiff modes.
    \item \textbf{Raw AI output:} A-stability was incorrectly treated as
    implying L-stability.
    \item \textbf{Risk audited:} The stability function satisfies
    $R_G(z)\to1$ as $z\to-\infty$, so infinitely stiff modes are not strongly
    damped.
    \item \textbf{Verification and action:} The report derives
    $R_G(z)=(1+z/2+z^2/12)/(1-z/2+z^2/12)$, proves A-stability, and explicitly
    records that Gauss IRK4 is not L-stable.
\end{itemize}

\subsection*{Changes Made to AI Output}
No numerical values were accepted from AI without reproduction. Numerical
claims were checked against the submitted scripts and data; analytical claims
were re-derived for the matrix problem; and generic wording was rewritten to
match the non-commutative matrix algebra and Lyapunov diagnostics used here.

\subsection*{Intellectual Ownership Declaration}
The mathematical formulations, numerical implementations, experiments,
validation evidence, figures, and final conclusions are the team's work.
Generative AI was used for drafting, clarification, and preliminary checking.
The authors take responsibility for the final mathematical and numerical
claims.
